\chapter{Weak-alpha source entry and the unit-\(n\) terminal}
\label{component:p3a3949869b9e}
\noindent\textbf{Component }\path{p3a3949869b9e}\par
{\small\noindent Source: \path{NEW/global/weak_alpha_source_and_unit_n_REDERIVED.tex}\par}
{\small\noindent SHA256: \path{3a3949869b9e658c08d8de681cc2d74a105dc146f4db0ae07a00ec905aef7e72}\par}
{\small\noindent Inventory use: WA1-WA6 general weak-alpha entry and WA1-WA27 complete unit-n weak-m C2 terminal | WA1-WA27 common canonical weak-alpha source entry plus complete unit-n weak-m C2 all-tests theorem | WA1-WA10 only on weak-n source domain; exact full-parent/scalar/contact/normal identities WA12-WA20\par}
\medskip
\section{Weak-alpha source entry and the unit-\(n\) terminal}
\label{p3a3949869b9e:sec:new-weak-alpha-unit-n}

\paragraph{Scope.}
All coefficients are bounded in the source field, \(d=1\), and the
original canonical source is physical with an actual finite noncommuting
zero contact. The C2 convention is \(\lambda=\beta=0\), and the C1
convention is \(\lambda=n=0\). The first subsection applies to either
family with \(\alpha\) infinitesimal. The all-tests theorem proved below
has the additional C2 hypotheses that \(n\) is a positive unit and \(m\)
is infinitesimal. It includes all independent testing states, signs and
literal projective endpoints. The weak-\(n\) and C1 terminals are separate.

\subsection{Source entry, proved before any shape quotient is used}

Physicality forces \(C,A,D\) to be positive units. Here are exact
residue-face tests proving this assertion. In C2 the negative central
coefficient at \(A=0\), and its vanishing positive Laurent terms along
\(U=t^2,V=1/t\), are the displayed forbidden-face proof in the actual
positive-\(d\) source partition. If also \(C=0\) in residue, the source
has only the four alternating events \(n,m,A,D\). At \(u=v=1\) the full
parent at state \((1,-\varepsilon,\varepsilon,1)\) is exactly
\[
 16\varepsilon\{(1+n+m+A+D)\varepsilon-1\}<0
 \tag{WA1}
\]
for a sufficiently small positive real \(\varepsilon\).
Thus \(C\) cannot be weak. The remaining central coefficient identities
are
\[
\begin{aligned}
h|_{\alpha=D=0}
 &=(2A+2C+1)\{8\beta mn-2\beta-4mn-2m-2n-1\},\\
h|_{\alpha=C=n=0}
 &=-(2A+1)(2D+1)(2\beta+2m+1),\\
h|_{\alpha=A=n=0}
 &=-(2C+2D+1)(2\beta+2m+1).
\end{aligned}\tag{WA2}
\]
In the first line \(n\beta=0\), so it is strictly negative. For C2
with \(\alpha=D=0\), the other Laurent terms are
\(q_{20}U/V+q_{21}U\), killed by \(U=t^2,V=t\).
For C1 with \(\alpha=D=0\), they are
\(q_{12}V+q_{21}U+q_{22}UV\), killed by \(U=V=t\).
For C1 with \(\alpha=C=0\), they are \(q_{02}V/U+q_{12}V\),
killed by \(U=1/t,V=t\). For C1 with \(\alpha=A=0\), they are
\(q_{02}V/U+q_{12}V+q_{22}UV\), killed by \(U=t,V=t^2\).
Here \(t\) tends to zero through positive reals. These contradictions
prove all the asserted unit bounds without a source-contact assumption.

Write \(U=u_c^2,V=v_c^2>0\). The actual determinant dictionary now has
\[
 q_{01}=\alpha m a_1,\quad q_{02}=ma_2,\quad
 q_{10}=\alpha a_3,\quad q_{12}=a_4,\quad
 q_{20}=na_5,\quad q_{21}=a_6,\quad q_{22}=\beta a_7,
 \tag{WA3}
\]
where the \(a_i\) are positive bounded units. The two angular equations
and their sum are
\[
\begin{aligned}
 U^2(na_5+a_6V+\beta a_7V^2)&=mV(\alpha a_1+a_2V),\\
 V^2(ma_2+a_4U+\beta a_7U^2)&=\alpha a_3U+na_5U^2,\\
 a_6U+a_4V+2\beta a_7UV
 &=\alpha(ma_1/U+a_3/V).
\end{aligned}\tag{WA4}
\]
Every \(t_i\) in the dictionary is a positive unit because it includes
\(d=1\). At contact all noncentral Laurent terms are nonnegative,
and their sum is bounded. In particular \(U,V,mV/U,nU/V\) are bounded.
Multiplying the last equation by \(V\) shows
\(\alpha/V=O(U+V)\). The first two equations then exclude either
\(U\) or \(V\) having nonzero residue: if \(U\) is a unit the first
forces \(V\) a unit, and conversely the second forces \(U\) a unit;
the last equation then contradicts weak \(\alpha\).
Both are infinitesimal. Consequently
\[
 \alpha/V=o(1),\qquad
 m\asymp (U/V)^2(n+V),\qquad
 \alpha\asymp V(U+V).
 \tag{WA5}
\]
For the last comparison, the denominator
\(a_3+ma_1V/U\) is a positive bounded unit. Actual finite contact
forces \(m,\alpha>0\) in WA4. If \(n\) is a unit and \(m\) weak,
\[
 U/V\asymp\sqrt m=o(1).
 \tag{WA6}
\]
If \(n\) is weak, then, whether \(U/V\) is bounded or unbounded,
WA5 gives \(mV/U,nU/V=o(1)\). These same two conclusions hold
under WA6. They are the common source bounds used in the shape argument.

Orient the actual signs positively, use \(W_c=1\), and put
\[
 p=v_c,\quad q=Z_c,\quad P=p^2,\quad Q=q^2,\qquad
 u_c=rpq,\quad X_c=xpq^2,\quad Y_c=ypq.
 \tag{WA7}
\]
The exact core sign identity SE9 ensures \(W_cZ_c>0\), so this
normalization is legitimate. To prove boundedness before using these
quotients, eliminate the \(X,Y\) block inside the source field. Set
\[
 H_X=n+m+\beta+\tfrac12,\quad K=C+D+\tfrac12,\quad H=K+A .
\]
The block has diagonal positive residue \(H_X,H\) and an infinitesimal
off-diagonal entry. The source Schur entries satisfy
\[
 S_{WW}/U\sim a=AK/H>0,\qquad
 -S_{WZ}/(u_cv_c)\sim b=\tfrac12+AD/H>0.
 \tag{WA8}
\]
Indeed the only additional leading contributions are multiples of
\(mV/U,nU/V,mn\), all infinitesimal on the domains just stated.
The kernel equation gives \(r\sim b/a\), a positive unit. The scaled
\(X,Y\) equations are then a bounded system with positive triangular
residue; its apparent \(m/Q\) term is bounded since
\(m/Q=r^2mV/U\). Thus \(x,y\) are bounded before taking residues.
Those equations and the two remaining kernel equations give
\[
 Ay\sim\tfrac12,\quad
 K(r-y)\sim D+\tfrac12,\quad
 D(1+y-r)\sim r/2,\qquad
 x\sim-\frac{m+\beta+1/2}{H_X}\,r.
 \tag{WA9}
\]
For the last kernel division the possible extra \(\alpha/V\),
\(nU/V\) and \(mV/U\) terms vanish by WA5. In particular \(y,r-y\)
and \(1+y-r\) are positive units, and the critical residue dictionary is
\[
 A_0=\frac1{2y_0},\qquad
 C_0=\frac{1+y_0}{2(r_0-y_0)},\qquad
 D_0=\frac{r_0}{2(1+y_0-r_0)}.
 \tag{WA10}
\]
In the unit-\(n\), weak-\(m\) domain, \(\beta=0\), and hence
\[
 P,Q=o(1),\qquad m\asymp Q^2,\qquad
 \alpha\asymp P^2,\qquad
 x_0=-r_0/(1+2n_0).
 \tag{WA11}
\]
No independent testing coordinate occurred in this entry proof.

\subsection{The exact scalar and every first normal}

For reference the full finite quarter-parent is
\[
\begin{aligned}
F={}&nD_uD_vX^2+m(vW-X+uvY-uZ)^2
       +CD_v(uW-Y)^2+AD_uD_vY^2\\
 &+(D+\tfrac12)(uW+uvX-Y-vZ)^2
       +(\alpha+\tfrac12)D_u(vY-Z)^2\\
 &+(\beta+\tfrac12)D_v(X+uZ)^2-\tfrac12D_uD_vZ^2 .
\end{aligned}\tag{WA12}
\]
Here \(D_u=1+u^2,D_v=1+v^2\). We retain \(\beta\) symbolically to make
the exact common graph explicit, but the current terminal has \(\beta=0\).
In the testing chart \(W=1,Z=z\), use
\[
 u_g=rvz,\qquad X_g=xvz^2,\qquad Y_g=yvz,\qquad
 V_t=v^2,\quad T=z^2.
 \tag{WA13}
\]
In the rest of this subsection write \(V=V_t\) only. The full graph
restriction is
\[
\begin{aligned}
\Phi={}&mV+\alpha T+hVT+\mathcal C V^2T+EVT^2
   +\mathcal F V^2T^2+GV^3T^2+JV^2T^3+\mathcal K V^3T^3,\\
\mathcal C={}&Ay^2+C(r-y)^2+(\alpha+\tfrac12)y^2+2mry,\\
E={}&\alpha r^2+(m+\beta+\tfrac12)(r+x)^2+nx^2,\\
G={}&r^2y^2(A+\alpha+m+\tfrac12),\qquad
J=nr^2x^2,\qquad \mathcal K=(D+n+\tfrac12)r^2x^2,\\
\mathcal F={}&Ar^2y^2+(2D+1)rx(r-y-1)-(2\alpha+1)r^2y\\
 &-2mry(r+x)+\beta(r+x)^2+nx^2+rx+x^2/2 .
\end{aligned}\tag{WA14}
\]
The coefficient \(h\) is defined by this entire expansion; it is bounded.
Every displayed coefficient is bounded. In the current unit-\(n\)
domain \(\mathcal C,E\) are positive units, since
\[
 \mathcal C_0=r_0(1+y_0)/2,\qquad
 E_0=n_0r_0^2/(1+2n_0).
 \tag{WA15}
\]
The coefficient \(\mathcal F\) need not be positive.

The actual contact gives \(\Phi(P,Q)=\Phi_V(P,Q)=\Phi_T(P,Q)=0\).
Solving these three linear equations in \(m,\alpha,h\) is exact:
\[
\begin{aligned}
m={}&EQ^2+\mathcal F PQ^2+GP^2Q^2+2JPQ^3+2\mathcal K P^2Q^3,\\
\alpha={}&\mathcal C P^2+\mathcal F P^2Q+2GP^3Q
                         +JP^2Q^2+2\mathcal K P^3Q^2,\\
h={}&-2\mathcal CP-2EQ-3\mathcal FPQ-4GP^2Q
                         -4JPQ^2-5\mathcal K P^2Q^2 .
\end{aligned}\tag{WA16}
\]
Put
\[
 R=z(V-P),\quad S=v(T-Q),\quad
 \Psi=QR^2+PS^2+vzRS .
\]
With the exact balanced monomials
\[
\begin{aligned}
B_{32}&=2PQR^2+(vzR+PS)^2,\\
B_{23}&=2PQS^2+(QR+vzS)^2,\\
B_{33}&=(2PQ^2+VT^2)R^2+(2P^2Q+PVT)S^2
                              +(3PQ+PT)vzRS,
\end{aligned}\tag{WA17}
\]
the whole scalar, not a leading truncation, is
\[
 \Phi=\mathcal CR^2+ES^2+\mathcal F\Psi
                          +GB_{32}+JB_{23}+\mathcal K B_{33}.
 \tag{WA18}
\]
For \(v,z\) infinitesimal, all entries after the first two positive
unit pivots are infinitesimal, by \(P,Q=o(1)\). This is a strict
two-row Gram even when \(\mathcal F\) is signed.

Let \(L_u,L_X,L_Y\) be half the full first derivatives of WA12 at
WA13. The complete polynomials have the forms
\[
\begin{aligned}
 L_u&=vz A_u(V,T),\qquad L_Y=vz A_y(V,T),\\
 A_u&=a_{00}+a_{10}V+a_{01}T+a_{11}VT
                      +a_{12}VT^2+a_{21}V^2T+a_{22}V^2T^2,\\
 A_y&=e_{00}+e_{10}V+e_{11}VT+e_{21}V^2T,\\
 L_X&=v[-m+TB_X(V,T)],\qquad
 B_X=c_0+c_1V+c_2VT+c_3V^2T .
\end{aligned}\tag{WA19}
\]
All coefficients are bounded, and the expansion fixes every one
unambiguously. In particular
\[
 c_2=nr^2x,\qquad c_3=(D+n+\tfrac12)r^2x .
\]
Their actual source jets are \(A_u(P,Q)=A_y(P,Q)=0\) and
\(B_X(P,Q)=m/Q\). Finite monic difference division gives
\[
\begin{aligned}
 L_i&=v A_{i,V}(V,T)R+z B_i(T)S,\qquad i=u,Y,\\
 L_X&=vz B_{X,V}(V,T)R+
                \{T(Pc_2+P^2c_3)+m/Q\}S .
\end{aligned}\tag{WA20}
\]
Here \(A_{i,V}=[A_i(V,T)-A_i(P,T)]/(V-P)\), and \(B_{X,V}\)
is defined identically. The other divided polynomials are
\[
\begin{aligned}
B_u(T)&=a_{01}+Pa_{11}+Pa_{12}(T+Q)+P^2a_{21}
                                      +P^2a_{22}(T+Q),\\
B_Y(T)&=Pe_{11}+P^2e_{21}.
\end{aligned}
\]
These are literal polynomial identities with bounded coefficients,
including \(v=0\) or \(z=0\). Since \(m/Q=O(Q)=o(1)\),
every full first normal is an infinitesimal linear combination of \(R,S\).

Set \(N_u=u-rvz,N_X=X-xvz^2,N_Y=Y-yvz\). In the chart under
consideration these are infinitesimal. The full higher-normal
remainder is a polynomial in the three normals with no terms of
normal degree below two; all terms of degree at least three give
infinitesimal Gram entries. Its quadratic residue is
\[
 K_0(N_u-N_Y)^2+A_0N_Y^2+(n_0+\tfrac12)N_X^2\succ0.
 \tag{WA21}
\]
WA18 and WA20 therefore give a strict bounded Gram for the entire
parent on the three hard and two soft rows. The strict-Gram lemma
proves its tensor-SOS assertion, with no mixed inverse.

\subsection{The complementary local states}

Normalize a maximal testing state coordinate. At \(u,v=o(1)\)
the state residue has positive pivots on \(X,Y\); if one is a unit,
the full parent has positive residue immediately. Otherwise one
of \(W,Z\) is a unit. The preceding proof covers \(W=1,Z=o(1)\).
Every remaining state can be normalized to \(Z=1,W=\omega\)
bounded. Retain the entire positive \(m,\alpha\) packets.
The remainder lies exactly in the square of \((u,v,X,Y)\).
Its quadratic real residue is
\[
\begin{aligned}
 &(n_0+\tfrac12)\left(X+\frac{u}{1+2n_0}\right)^2
     +\frac{n_0}{1+2n_0}u^2\\
 &+H_0\left(Y-\frac{K_0\omega u-D_0v}{H_0}\right)^2
     +a_0(\omega u-r_0v)^2,\qquad a_0=A_0K_0/H_0>0 .
\end{aligned}\tag{WA22}
\]
The critical identities WA10 prove the equality. This quadratic
form is positive definite for every real \(\omega\), including zero.
All higher normal terms are infinitesimal corrections on these
four rows, while \(\omega-\bar\omega\) is retained in coefficients.
Thus every state at the double-infinitesimal angle corner is covered.

\subsection{All other testing angles, including the literal endpoints}

Again retain the entire \(m,\alpha\) packets and call the remainder
\(F_0\). Its source residue is the critical \(C,A,D,n\) source.
For bounded \(u,v\), \(F_0\) is exactly a bounded quadratic form
in the four rows \((uW,X,Y,Z)\). Its real \(Z=0\) block is positive:
the \(nX^2\) and \(AY^2\) packets and \(C(uW-Y)^2\) suffice.
The determinant is
\[
 \det \operatorname{Mat}_{(uW,X,Y,Z)}(\bar F_0)
 =D_u^2D_v^2\{q_{20}u^2+q_{12}v^4+q_{21}u^2v^2\}/16>0
 \tag{WA23}
\]
whenever at least one of the two real angle residues is nonzero.
Here \(q_{20},q_{12},q_{21}>0\) are the actual residue determinant
coefficients, and \(h=0\) follows from WA10. This proves strictness
of the full four-row Gram. Literal \(u=0\) is allowed.

For Bob reciprocal \(b=1/v\) infinitesimal and Alice bounded,
use the same rows \((uW,X,Y,Z)\) in \(b^2F_0(u,1/b)\).
At \(b=0\) its whole real expression is
\[
 C_0(uW-Y)^2+(A_0+\tfrac12)D_uY^2
             +D_0(uX-Z)^2+(n_0+\tfrac12)D_uX^2\succ0 .
 \tag{WA24}
\]
It is strict at every real \(u\), including zero.
For Alice reciprocal \(s=1/u\) infinitesimal and Bob bounded,
the original state rows suffice. The expression at \(s=0\) is
\[
 n_0D_vX^2+C_0D_vW^2+A_0D_vY^2
          +(D_0+\tfrac12)(W+vX)^2+\tfrac12(vY-Z)^2\succ0 .
 \tag{WA25}
\]
This includes \(v=0\), using the positive \(n_0\) pivot.

For both reciprocal angles \(s,b=o(1)\), use the five rows
\((W,X,Y,sZ,bZ)\). The whole reciprocal polynomial belongs
to their square ideal, and its real coefficient Gram is
\[
 C_0W^2+A_0Y^2+D_0(X-sZ)^2
          +(n_0+\tfrac12)X^2+\tfrac12(Y-bZ)^2\succ0 .
 \tag{WA26}
\]
Every remaining coefficient is infinitesimal. These exact finite
square-ideal allocations include \(s=0\), \(b=0\) and arbitrary
states; no projective endpoint is obtained by a limiting argument.

The four angle charts and the two local state charts are exhaustive.
Testing sign gauges, positive pure-field homogeneous normalization
factors and real/imaginary state splitting restore the original parent:
\[
 F_c(u,v;\psi)\ \in\ \Sigma(K_c\otimes_{\mathbb R}K_t)^2
 \quad\text{for every independent projective test and state.}
 \tag{WA27}
\]
This closes exactly the C2 unit-\(n\), weak-\(m,\alpha\) terminal.
The common source-entry subsection does not by itself close the
weak-\(n\) or C1 weak-\(\alpha\) branches.

\paragraph{Evidence boundary.}
The focused checker expands the full parent, the scalar and all first
normals, checks the balanced identities with their explicit contact
remainders, and verifies the complete hard and outer residue forms.
The positive-unit source estimates, the finite quantified chart cover
and the application of the separated strict-Gram lemma are the
analytical proof above. No unavailable historical local proof is used.
